\clearpage\pagestyle{otMain}\fontsize{11}{15}\selectfont\corrigesChapitre{4}{Développer, factoriser et prouver}
\begin{blocLivre}{cor-4-1}\begin{corrige}{1}{Remplacer une lettre}\demarche{Écrire la substitution avant de simplifier.}
Pour $x=-2$, $A=2\times4+6+1=\res{15}$. Pour $x=3$, $A=18-9+1=\res{10}$. Les parenthèses autour de $-2$ évitent de calculer un carré avec le mauvais signe.
\end{corrige}\end{blocLivre}
\begin{blocLivre}{cor-4-2}\begin{corrige}{2}{Réduire une expression}\demarche{Regrouper seulement les termes de même nature.}
Les termes en $x$ donnent $(7+3-2)x=8x$. Les constantes donnent $-4+9=5$. On obtient $\res{8x+5}$. Les deux sortes de termes restent distinctes.
\end{corrige}\end{blocLivre}
\begin{blocLivre}{cor-4-3}\begin{corrige}{3}{Prendre l’opposé}\demarche{Distribuer chaque facteur à toute sa parenthèse.}
Le signe moins change les deux signes : $-4x+7+6-2x=\res{-6x+13}$. Le produit par 2 porte sur les deux termes de la deuxième parenthèse.
\end{corrige}\end{blocLivre}
\begin{blocLivre}{cor-4-4}\begin{corrige}{4}{Double distributivité}\demarche{Écrire les quatre produits avant la réduction.}
On écrit les quatre produits : $3x^2-15x+2x-10$. Les deux termes en $x$ se regroupent, d’où $\res{3x^2-13x-10}$.
\end{corrige}\end{blocLivre}
\begin{blocLivre}{cor-4-5}\begin{corrige}{5}{Un facteur commun}\demarche{Rechercher un facteur numérique et éventuellement littéral.}
Le premier facteur commun est 7 : $\res{7(2x+3)}$. Pour le second, 4 et $x$ sont communs : $\res{4x(3x-2)}$. Redévelopper retrouve les expressions initiales.
\end{corrige}\end{blocLivre}
\begin{blocLivre}{cor-4-6}\begin{corrige}{6}{Différences de carrés}\demarche{Identifier les deux carrés avant d’appliquer l’identité.}
On reconnaît $81=9^2$ et $25x^2=(5x)^2$. Les résultats sont $\res{(x-9)(x+9)}$ et $\res{(5x-2)(5x+2)}$.
\end{corrige}\end{blocLivre}
\begin{blocLivre}{cor-4-7}\begin{corrige}{7}{Deux carrés de sommes}\demarche{Conserver le double produit.}
Le premier développement est $\res{x^2+12x+36}$. Le second est $4x^2-12x+9$. Dans chaque cas, le terme du milieu est le double produit des deux termes de la parenthèse, avec son signe.
\end{corrige}\end{blocLivre}
\begin{blocLivre}{cor-4-8}\begin{corrige}{8}{Calculer astucieusement}\demarche{Choisir un nombre rond comme point d’appui.}
$103\times97=(100+3)(100-3)=10000-9=\res{9991}$. Ensuite $(50-1)^2=2500-100+1=\res{2401}$.
\end{corrige}\end{blocLivre}
\begin{blocLivre}{cor-4-9}\begin{corrige}{9}{Factoriser en deux temps}\demarche{Repérer un facteur commun avant une identité éventuelle.}
Dans la première expression, la parenthèse entière est commune : $\res{(x-2)(3x+5)}$. Dans la seconde, sortir 9 puis utiliser une différence de carrés : $\res{9(x-2)(x+2)}$.
\end{corrige}\end{blocLivre}
\begin{blocLivre}{cor-4-10}\begin{corrige}{10}{Tester une conjecture}\demarche{Utiliser un contre-exemple puis expliquer l’erreur algébrique.}
Pour $x=1$, $9-1=8$ : l’affirmation est fausse. Le développement donne $(x^2+4x+4)-x^2=\res{4x+4}$. Le résultat vaut 4 uniquement pour $x=0$.
\end{corrige}\end{blocLivre}
\begin{blocLivre}{cor-4-11}\begin{corrige}{11}{Deux programmes équivalents}\demarche{Traduire les étapes avec un nombre quelconque.}
Pour un nombre $x$, A donne $4(x+5)-8=4x+20-8=4x+12$. B donne directement $4x+12$. Les deux expressions sont identiques pour tout $x$. Tester seulement quelques nombres n’aurait pas suffi.
\end{corrige}\end{blocLivre}
\begin{blocLivre}{cor-4-12}\begin{corrige}{12}{Une bordure}\demarche{Soustraire les aires en calculant d’abord le bon côté extérieur.}
Le grand carré a pour côté $x+4$, puisque la bordure est des deux côtés. L’aire cherchée vaut $(x+4)^2-x^2=\res{8x+16}$ cm$^2$. Pour $x=7$, elle vaut $\res{72\text{ cm}^2}$. On contrôle $11^2-7^2=72$.
\label{fig-cor-4-12}\figureGeo[corrige]{bordure}
\end{corrige}\end{blocLivre}
\begin{blocLivre}{cor-4-13}\begin{corrige}{13}{Des entiers consécutifs}\demarche{Écrire la forme générale de chaque parité.}
Un impair s’écrit $2n+1$ avec $n$ entier. Son carré est $4n^2+4n+1=2(2n^2+2n)+1$, donc impair. Un pair s’écrit $2n$ ; son carré vaut $4n^2$, donc est bien multiple de 4. Les preuves portent sur un entier quelconque.
\end{corrige}\end{blocLivre}
\begin{blocLivre}{cor-4-14}\begin{corrige}{14}{Une égalité impossible ?}\demarche{Confronter la solution algébrique au domaine géométrique.}
La différence vaut $(x+3)^2-x^2=6x+9$. Pour obtenir 9, il faudrait $x=0$, exclu par la positivité du côté. Pour 39, $6x+9=39$ donne $x=5$, qui convient. Les côtés sont alors 8 et 5 ; $64-25=39$.
\end{corrige}\end{blocLivre}
\begin{blocLivre}{cor-4-15}\begin{corrige}{15}{Une différence toujours positive}\demarche{Comparer les deux quantités par leur différence.}
On calcule la différence : $A-B=x^2+2x+1-(x^2+2x)=1$. Donc $\res{A=B+1}$ pour tout $x$ ; A est toujours strictement supérieur à B. Cette méthode évite de distinguer les signes des facteurs de B.
\end{corrige}\end{blocLivre}
